% Values: PRECISE, ECONOMICAL, FAITHFUL
% Evidence planned:
% - PRECISE: a table row for every derived form of the book plan's "Derived
%   Forms"; the grouping of a let and of a module body stated with displayed
%   definitions (bound variables, right-hand sides, "refers to", components,
%   numbering); every equivalent form closed under B4; fresh identifiers
%   stated (none).
% - ECONOMICAL: if, else-if chains, &&, ||, negation, lists, the field
%   accessor, field access, list patterns, lambdas with several patterns, function
%   definitions, let with several definitions and the order of module
%   declarations carry no rules of their own.
% - FAITHFUL: A56 (infixed && and || are conditionals, the operator values
%   are strict); A16, A30 as amended (every module has a header, so no
%   derived form supplies one); A57, A58 (order among independent components is not
%   significant); A26 (restrictions apply to definitions as written); A4
%   (field access); D5 (accessor and access through record patterns).
\chapter{Derived Forms}\label{ch:appendix-derived-forms}

Several derived forms are provided in the Core and in Modules; they are
presented in Figures~\ref{fig:appendix-derived-forms-expressions},
\ref{fig:appendix-derived-forms-patterns},
\ref{fig:appendix-derived-forms-bindings}
and~\ref{fig:appendix-derived-forms-modules}. Each derived form is given with
its equivalent form, and each row of the tables is a rewriting rule
\[
  \text{Derived form} \;\Longrightarrow\; \text{Equivalent form}
\]
which may be applied to any subphrase, repeatedly, until the phrase is a
phrase of the bare language of Chapters~\ref{ch:syntax-core}
and~\ref{ch:syntax-modules}. Within a figure the rows are ordered, and the
figures follow one another, so that every equivalent form uses only bare
syntax and the derived forms of earlier rows. % B4
Two rows, the grouping of the definitions of a \kw{let} and the grouping of the
declarations of a module, are not local rewritings; their equivalent forms
are defined in Sections~\ref{sec:appendix-derived-forms-let}
and~\ref{sec:appendix-derived-forms-modules}. See
Appendix~\ref{ch:appendix-full-grammar} for the full grammar, including all
the derived forms.

In equivalent forms, \kw{Basics.True} and \kw{Basics.False} are the
constructors of the type $\tBool$, \kw{Basics.negate} is the negation
function, and \kw{::} is the list constructor. They are looked up as any
other identifier of the module in which the derived form occurs. Because
only an import of the module \kw{Basics} provides the qualifier \kw{Basics}
(Section~\ref{sec:syntax-modules-restrictions}), \kw{Basics.True},
\kw{Basics.False} and \kw{Basics.negate} always denote the bindings of the
initial basis (Appendices~\ref{ch:appendix-static-basis}
and~\ref{ch:appendix-dynamic-basis}), whatever the module imports or
declares; \kw{::} is bound in the initial environment $E_0$ of those
appendices, from which every module is elaborated and evaluated, and no
program declares it. Being qualified, or an operator, none of them is a free
occurrence of a variable in the sense below.
% Audit X18, V3: resolved by the restriction on the qualifier Basics (A89),
% not by the provenance of the identifier; the operator :: is fixed by the
% operator table of the initial basis (design decision 1).

In the first row of Figure~\ref{fig:appendix-derived-forms-expressions},
$\expm_1\ \opv\ \expm_2$ is an infixed application as grouped by the
resolution of Section~\ref{sec:syntax-core-infix}: $\expm_1$ and $\expm_2$
are the two operands of $\opv$ after grouping, and $\opv$ is neither \kw{\&\&}
nor \kw{||}. Each operator of a chain is rewritten in this way, outermost
first. The infixed forms of \kw{\&\&} and \kw{||} have rows of their own, which
make them conditionals; the operator values \kw{(\&\&)} and \kw{(||)} are
bare, and denote functions of the initial basis that take both arguments as
values. % A56
In the row for negation, $\kw{-}\atexp$ is a negation as defined in
Section~\ref{sec:syntax-core-lexical}; it is atomic, so that it may be the
argument of an application. % DECIDED: negation is placed among the atomic
% expressions (the survey: in `f -x` the negated term is an argument); its
% equivalent form is parenthesised, so it is correct at either level.
In the row for field access, the atomic expressions that may precede
$\kw{.}\lab$ are restricted in Section~\ref{sec:syntax-core-restrictions}
(the restrictions on records); after any other atomic expression, and after a
field accessor, a \kw{.} begins a field accessor
(Section~\ref{sec:syntax-core-lexical}), so $\kw{.}\lab_1\kw{.}\lab_2$ is the
application of $\kw{.}\lab_1$ to $\kw{.}\lab_2$.
% A4
In the equivalent form of the field accessor $\kw{.}\lab$, the record pattern
$\kw{\kwlb}\ \lab\ \kw{\kwrb}$ binds the variable spelt as the label $\lab$,
and the body of the lambda is that variable. % D5

An annotation $\var\ \kw{:}\ \ty$ and the definition of $\var$ that follows it
(Chapter~\ref{ch:syntax-core}) are written together, here as in the bare
grammar, as one value binding with $\kw{:}\ \ty$ after $\var$. Thus a
definition is a value binding without \kw{and}, of the form
$\var\ \opt{\kw{:}\ \ty}\ \atpat_1\cdots\atpat_n\ \kw{=}\ \expm$ ($n\geq 0$)
or $\atpat\ \kw{=}\ \expm$.

The syntactic restrictions of Chapters~\ref{ch:syntax-core}
and~\ref{ch:syntax-modules} apply to a program as written, before any derived
form in it is replaced by its equivalent form
(Section~\ref{sec:syntax-core-restrictions}), and are not checked again on
equivalent forms. In
particular the variable bound in the equivalent form of $\kw{.}\lab$ may have
the name of a variable in scope, and a restriction that distinguishes a
definition with parameters from one without refers to the definition as
written. % A26
% DECIDED: restrictions are checked on the phrase as written. Otherwise the
% equivalent form of `.x` would break the no-shadowing restriction when a
% variable x is in scope, and `f n = f n` and `f = \n -> f n` (which A26
% treats differently in a let) could not be told apart after rewriting.

An occurrence of a variable in a phrase is \emph{free} if it is unqualified
and is not within the scope of a binding of that variable, by a pattern or by
a value binding, made inside the phrase. No equivalent form adds or removes a
free occurrence of a variable, so the groupings of
Sections~\ref{sec:appendix-derived-forms-let}
and~\ref{sec:appendix-derived-forms-modules} are the same whichever other
derived forms have already been rewritten.

No equivalent form introduces a new identifier, so no table carries a note of
fresh identifiers. % B5
The only variable bound by an equivalent form is the one bound by the
equivalent form of $\kw{.}\lab$, which is spelt as the label.

\begin{figure}[htbp]
\centering
\par\noindent\textbf{Expressions}\ \expm\par\smallskip
\begin{derivedforms}
\dform{\citelabel{df:infix}\expm_1\ \opv\ \expm_2}
      {\kw{(}\opv\kw{)}\ \kw{(}\expm_1\kw{)}\ \kw{(}\expm_2\kw{)}}
\hline
\dform{\citelabel{df:if}\kw{if}\ \expm_1\ \kw{then}\ \expm_2\ \kw{else}\ \expm_3}
      {\begin{array}[t]{@{}l@{}}
         \kw{case}\ \expm_1\ \kw{of}\ \kw{Basics.True}\ \kw{->}\ \kw{(}\expm_2\kw{)}\\
         \phantom{\kw{case}\ \expm_1\ \kw{of}}\ \kw{;}\ \kw{Basics.False}\ \kw{->}\ \expm_3
       \end{array}}
\hline
\dform{\citelabel{df:if-chain}\begin{array}[t]{@{}l@{}}
         \kw{if}\ \expm_1\ \kw{then}\ \expm'_1\\
         \kw{else}\ \kw{if}\ \expm_2\ \kw{then}\ \expm'_2\ \cdots\\
         \kw{else}\ \kw{if}\ \expm_n\ \kw{then}\ \expm'_n\\
         \kw{else}\ \expm
       \end{array}}
      {\begin{array}[t]{@{}l@{}}
         \kw{if}\ \expm_1\ \kw{then}\ \expm'_1\ \kw{else}\ \kw{(}\\
         \;\kw{if}\ \expm_2\ \kw{then}\ \expm'_2\ \kw{else}\ \kw{(}\ \cdots\\
         \;\;\kw{if}\ \expm_n\ \kw{then}\ \expm'_n\ \kw{else}\ \expm
           \ \kw{)}\cdots\kw{)}\\
         \qquad (n\geq 2)
       \end{array}}
\hline
\dform{\citelabel{df:and}\expm_1\ \kw{\&\&}\ \expm_2}
      {\kw{if}\ \expm_1\ \kw{then}\ \expm_2\ \kw{else}\ \kw{Basics.False}}
\hline
\dform{\citelabel{df:or}\expm_1\ \kw{||}\ \expm_2}
      {\kw{if}\ \expm_1\ \kw{then}\ \kw{Basics.True}\ \kw{else}\ \expm_2}
\hline
\dform{\citelabel{df:lambda-params}\kw{\textbackslash}\atpat_1\cdots\atpat_n\ \kw{->}\ \expm}
      {\kw{\textbackslash}\atpat_1\ \kw{->}\ \cdots\
       \kw{\textbackslash}\atpat_n\ \kw{->}\ \expm
       \quad (n\geq 2)}
\end{derivedforms}

\bigskip
\par\noindent\textbf{Atomic Expressions}\ \atexp\par\smallskip
\begin{derivedforms}
\dform{\citelabel{df:negation}\kw{-}\atexp}
      {\kw{(}\kw{Basics.negate}\ \atexp\kw{)}}
\hline
\dform{\citelabel{df:list}\kw{[}\expm_1\kw{,}\ \cdots\kw{,}\ \expm_n\kw{]}}
      {\kw{(}\expm_1\ \kw{::}\ \cdots\ \kw{::}\ \expm_n\ \kw{::}\ \kw{[ ]}\kw{)}
       \qquad (n\geq 1)}
\hline
\dform{\citelabel{df:field-accessor}\kw{.}\lab}
      {\kw{(}\kw{\textbackslash}\kw{\kwlb}\ \lab\ \kw{\kwrb}\ \kw{->}\ \lab\kw{)}}
\hline
\dform{\citelabel{df:field-selection}\atexp\kw{.}\lab}
      {\kw{(}\kw{.}\lab\ \atexp\kw{)}}
\end{derivedforms}
\caption{Derived forms of Expressions and Atomic Expressions}
\label{fig:appendix-derived-forms-expressions}
\end{figure}

\begin{figure}[htbp]
\centering
\par\noindent\textbf{Atomic Patterns}\ \atpat\par\smallskip
\begin{derivedforms}
\dform{\citelabel{df:list-pattern}\kw{[}\pat_1\kw{,}\ \cdots\kw{,}\ \pat_n\kw{]}}
      {\kw{(}\pat_1\ \kw{::}\ \kw{(}\pat_2\ \kw{::}\ \cdots\
       \kw{(}\pat_n\ \kw{::}\ \kw{[ ]}\kw{)}\cdots\kw{)}\kw{)}
       \quad (n\geq 1)}
\end{derivedforms}
\caption{Derived forms of Patterns}
\label{fig:appendix-derived-forms-patterns}
\end{figure}

In the equivalent form for lists, each \kw{::} is an infixed application
whose right operand is everything to its right:
\[
  \expm_1\ \kw{::}\ \kw{(}\expm_2\ \kw{::}\ \cdots\
  \kw{(}\expm_n\ \kw{::}\ \kw{[ ]}\kw{)}\cdots\kw{)}.
\]
In the equivalent form for list patterns, each \kw{::} is the bare pattern
$\pat\ \kw{::}\ \pat$, the parentheses make the result an atomic pattern,
and for $n=1$ the equivalent form is $\kw{(}\pat_1\ \kw{::}\ \kw{[ ]}\kw{)}$.
% Audit Y8: the list-pattern row is given directly, without recursion (B4).

\begin{figure}[htbp]
\centering
\par\noindent\textbf{Value Bindings}\ \valbind\par\smallskip
\begin{derivedforms}
\dform{\citelabel{df:function-binding}\begin{array}[t]{@{}l@{}}
         \var\ \opt{\kw{:}\ \ty}\ \atpat_1\cdots\atpat_n\ \kw{=}\ \expm\\
         \quad\opt{\kw{and}\ \valbind}
       \end{array}}
      {\begin{array}[t]{@{}l@{}}
         \var\ \opt{\kw{:}\ \ty}\ \kw{=}\
           \kw{\textbackslash}\atpat_1\cdots\atpat_n\ \kw{->}\ \expm\\
         \quad\opt{\kw{and}\ \valbind}
         \qquad (n\geq 1)
       \end{array}}
\end{derivedforms}

\bigskip
\par\noindent\textbf{Declarations}\ \dec\par\smallskip
\begin{derivedforms}
\dform{\citelabel{df:let-grouping}\kw{let}\ \valbind_1\ \cdots\ \valbind_n\ \kw{in}\ \expm}
      {\begin{array}[t]{@{}l@{}}
         \kw{let}\ \dec_1\ \kw{in}\ \cdots\ \kw{let}\ \dec_k\ \kw{in}\ \expm\\
         \text{(see Section~\ref{sec:appendix-derived-forms-let})}
         \qquad (n\geq 2)
       \end{array}}
\end{derivedforms}
\caption{Derived forms of Value Bindings and Declarations}
\label{fig:appendix-derived-forms-bindings}
\end{figure}

\begin{figure}[htbp]
\centering
\par\noindent\textbf{Top-level Declarations}\ \topdec\par\smallskip
\begin{derivedforms}
\dform{\citelabel{df:module-grouping}\topdec_1\ \cdots\ \topdec_n}
      {\begin{array}[t]{@{}l@{}}
         \kw{type}\ \datbind\ \kw{;}\\
         \kw{type alias}\ \typbind'_1\ \kw{;}\ \cdots\
           \kw{type alias}\ \typbind'_l\ \kw{;}\\
         \topdec'_1\ \kw{;}\ \cdots\ \kw{;}\ \topdec'_k\\
         \text{(see Section~\ref{sec:appendix-derived-forms-modules})}
         \qquad (n\geq 2)
       \end{array}}
\end{derivedforms}
\caption{Derived forms of Top-level Declarations}
\label{fig:appendix-derived-forms-modules}
\end{figure}

No derived form supplies the header of a module: every module is written
with one (Chapter~\ref{ch:syntax-modules}). % A16, A30

\section{Grouping the Definitions of a Let}
\label{sec:appendix-derived-forms-let}

A \kw{let} expression as written has definitions $\valbind_1,\ldots,\valbind_n$
($n\geq 2$), each a value binding without \kw{and}. The bare language has
$\kw{let}\ \dec\ \kw{in}\ \expm$ only for a single group of bindings joined by
the abstract \kw{and}, so the definitions are grouped and ordered as follows.
% Survey "Rewrites": let block; D7

The variables \emph{bound by} a definition, and its \emph{right-hand side},
are given by the following table, in which $m\geq 1$ and a definition of the
first two forms may carry an annotation $\kw{:}\ \ty$ after $\var$:
\[
  \begin{array}{lll}
    \text{Definition} & \text{Bound} & \text{Right-hand side}\\[2pt]
    \var\ \kw{=}\ \expm
      & \var & \expm\\
    \var\ \atpat_1\cdots\atpat_m\ \kw{=}\ \expm
      & \var & \kw{\textbackslash}\atpat_1\cdots\atpat_m\ \kw{->}\ \expm\\
    \atpat\ \kw{=}\ \expm
      & \text{those of}\ \atpat & \expm
  \end{array}
\]
The variables bound by a pattern are its variable patterns, the variables
spelt as the labels of its record patterns, and the variables following \kw{as}
in it.
A definition $\valbind_i$ \emph{refers to} a definition $\valbind_j$ if a
variable bound by $\valbind_j$ occurs free in the right-hand side of
$\valbind_i$; $j=i$ is allowed. Two definitions are in the same
\emph{component} if each can be reached from the other by a chain of zero or
more steps of ``refers to''. The components partition the definitions. A
component is \emph{recursive} if it has more than one member, or if its one
member refers to itself. The members of the recursive components are exactly
the definitions that lie on a cycle of ``refers to'', and the restriction on
recursion in Section~\ref{sec:syntax-core-restrictions} requires each of them
to be a definition with parameters; in particular a definition
$\atpat\ \kw{=}\ \expm$ is never in a recursive component. % A26

Number the components $1,\ldots,k$ so that whenever a member of component $a$
refers to a member of component $b\neq a$, then $b<a$. Such a numbering
exists, because a chain of ``refers to'' that returns to its start never
leaves a component. For each $a$, let
$\valbind_{i_1},\ldots,\valbind_{i_p}$, with $i_1<\cdots<i_p$, be the members
of component $a$, and let
\[
  \dec_a \;=\; \valbind_{i_1}\ \kw{and}\ \cdots\ \kw{and}\ \valbind_{i_p}.
\]
Then the equivalent form is given by
\[
  \begin{array}{l}
    \kw{let}\ \valbind_1\ \cdots\ \valbind_n\ \kw{in}\ \expm\\
    \quad\Longrightarrow\;
    \kw{let}\ \dec_1\ \kw{in}\ \kw{let}\ \dec_2\ \kw{in}\ \cdots\
    \kw{let}\ \dec_k\ \kw{in}\ \expm .
  \end{array}
\]
An annotation is part of the definition it annotates, and so goes with that
definition into $\dec_a$.

When two components are independent (neither can be reached from the other)
the numbering is not unique, and which admissible numbering is chosen does not
matter: independent components bind distinct variables and neither refers to
the other, the bindings of a component add no free type variable to the
context ($\Clos$ generalises every type variable not free in the context,
Section~\ref{sec:static-core-closure}), and evaluation has neither state nor
exceptions, so each component elaborates and evaluates to the same result
whichever independent components precede it. % A57, A58

\section{Grouping the Declarations of a Module}
\label{sec:appendix-derived-forms-modules}

The declarations of a module body as written stand in any order. Each
declaration $\topdec_i$ is one of the following, where $m\geq 0$:
\[
  \begin{array}{ll}
    \kw{type}\ \datbind_i & \text{a custom type, declaring one type;}\\
    \kw{type alias}\ \typbind_i & \text{an alias, declaring one alias;}\\
    \kw{port}\ \var\ \kw{:}\ \ty & \text{a port;}\\
    \var\ \opt{\kw{:}\ \ty}\ \atpat_1\cdots\atpat_m\ \kw{=}\ \expm
      & \text{a definition.}
  \end{array}
\] A module body has no definitions of the form
$\atpat\ \kw{=}\ \expm$ (Chapter~\ref{ch:syntax-modules}). A module body as
written is always rewritten by this row, even when its declarations already
stand in an order the row could produce. The equivalent form
\[
  \begin{array}{l}
    \kw{type}\ \datbind\ \kw{;}\\
    \kw{type alias}\ \typbind'_1\ \kw{;}\ \cdots\
      \kw{type alias}\ \typbind'_l\ \kw{;}\\
    \topdec'_1\ \kw{;}\ \cdots\ \kw{;}\ \topdec'_k
  \end{array}
\]
is built as follows, where \kw{;} is the abstract separator of top-level
declarations in the bare grammar (Chapter~\ref{ch:syntax-modules}); a part
that is omitted below is omitted with its separator.
% Survey "Rewrites": top-level values; A59

\begin{enumerate}
\item \emph{Custom types.} If $\kw{type}\ \datbind_{i_1},\ldots,
  \kw{type}\ \datbind_{i_p}$, with $i_1<\cdots<i_p$, are the custom types of
  the body, then all of them form one group,
  \[
    \datbind \;=\; \datbind_{i_1}\ \kw{and}\ \cdots\ \kw{and}\ \datbind_{i_p},
  \]
  whose members may refer to one another. If $p=0$, the declaration
  $\kw{type}\ \datbind$ is omitted.

\item \emph{Aliases.} An alias \emph{mentions} another alias of the body if
  the type constructor that the other binds occurs, unqualified, in its body.
  The aliases $\typbind'_1,\ldots,\typbind'_l$ are those of the body, ordered
  so that each comes after every alias it mentions. Such an order exists,
  because the restrictions of Chapter~\ref{ch:syntax-modules} forbid a chain
  of mentions that returns to its start. If $l=0$, there are no alias
  declarations.

\item \emph{Ports and definitions.} A port $\kw{port}\ \var\ \kw{:}\ \ty$
  binds $\var$ and has no right-hand side, so it refers to nothing; a
  definition binds its $\var$ and has the right-hand side of
  Section~\ref{sec:appendix-derived-forms-let}. Over the ports and
  definitions of the body, ``refers to'', components and their numbering
  $1,\ldots,k$ are defined as in Section~\ref{sec:appendix-derived-forms-let}.
  A port is alone in its component. For each $a$, $\topdec'_a$ is the port if
  component $a$ is a port, and otherwise is the value group
  \[
    \valbind_{i_1}\ \kw{and}\ \cdots\ \kw{and}\ \valbind_{i_q},
  \]
  where $\valbind_{i_1},\ldots,\valbind_{i_q}$, with $i_1<\cdots<i_q$, are the
  members of component $a$, each with its annotation.
\end{enumerate}

The type constructors bound by $\datbind$ and by the aliases are in scope
throughout these type declarations: a constructor binding in $\datbind$ may
mention an alias of the body, and an alias may mention a custom type of
$\datbind$. Chapter~\ref{ch:static-modules} elaborates the group $\datbind$
together with the aliases that follow it.
% DECIDED: a custom type may mention a local alias and an alias may mention a
% local custom type (survey, "Custom Types and Aliases": unions are entered
% before aliases, and constructors see all local types), so neither order of
% the group and the aliases is a sequence of independent declarations.
% Expanding the aliases into datbind, as SML's withtype does, is not a
% syntactic translation here: an alias of an extensible record applied to an
% imported alias has no expansion in the type syntax. The type declarations
% are therefore one unit for elaboration (a question for static-modules).

The values that the type declarations bind (the constructors of $\datbind$
and the constructors of record aliases) are thus in scope in every port and
definition. As for a \kw{let}, which order is chosen among aliases neither of
which mentions the other, and which admissible numbering is chosen for the
components, does not matter: an alias denotes the same type function
whichever unrelated aliases precede it, and a component of ports and
definitions elaborates and evaluates to the same result whichever independent
components precede it, for the reasons given in
Section~\ref{sec:appendix-derived-forms-let}. % A57, A58, A59
